% Version: 2026.1.0.0
% WHAT: Formal companion for artifact-read scope in Engineering Studio AI.
% WHY: Runtime artifact lookup needs a checkable, fail-closed read boundary.
% HOW: Prove an exact-slot invariant under explicit namespace assumptions.
\documentclass[12pt]{article}
\usepackage[utf8]{inputenc}
\usepackage[T1]{fontenc}
\usepackage{times}
\usepackage[margin=1in]{geometry}
\usepackage{amsmath,amssymb,amsthm}
\usepackage{booktabs,tabularx,ragged2e}
\usepackage{enumitem}
\usepackage[final]{microtype}
\usepackage{fancyhdr,lastpage}
\usepackage{xurl}
\usepackage[hidelinks]{hyperref}
\hypersetup{
  pdftitle={Exact-Slot Artifact Reads in Engineering Studio AI},
  pdfsubject={Engineering Studio AI formal companion}
}
\pagestyle{fancy}
\fancyhf{}
\lhead{Engineering Studio AI}
\rhead{Exact-slot artifact reads}
\cfoot{\thepage\ of \pageref{LastPage}}
\setlength{\headheight}{15pt}
\setlength{\emergencystretch}{2em}
\setlist[enumerate]{leftmargin=*}
\newtheorem{theorem}{Theorem}
\newtheorem{definition}[theorem]{Definition}
\newcommand{\denied}{\bot}
\newcommand{\regular}{\operatorname{regular}}
\newcommand{\slot}{\operatorname{slot}}
\newcommand{\resolve}{\operatorname{resolve}}
\newcommand{\select}{\operatorname{select}}

\title{Exact-Slot Artifact Reads in Engineering Studio AI\\
\large A formal companion to the OpenAI Dev Week paper}
\date{2026-11-20}

\begin{document}
\maketitle
\begin{abstract}
Engineering Studio exposes stage artifacts through JSON previews, single
downloads, and run-level archives. These surfaces rely on a shared lookup
that, at the examined baseline, checks run membership and file existence.
Existence alone does not establish which declared slot a selector may
read. This companion defines an exact-slot acceptance predicate and proves
its namespace safety, local completeness, and finite trace composition.
The model distinguishes grammar, registration, stage membership, canonical
resolution, and regular-file observation. It assumes a trusted storage
root and stable selection through open. Its conclusion is conditional
path confinement. Multi-user authorization, physical file identity
across names, content integrity, and attacker-controlled mutation require
separate controls and evidence.
\end{abstract}

\section*{Key Words}
Artifact isolation; namespace; canonical path; read boundary; audit;
stable observation; proof of refinement.

\section*{Executive Summary}
A selectable name must first match the declared identifier grammar and the
corresponding run and stage registries. Derive exactly one expected slot
from validated components and the canonical configuration root. Admit a
resolved path only if it equals that slot under the namespace comparison
contract and a regular-file check succeeds. The transport keeps its
existing missing-artifact behavior. Evaluation of an uploaded report or of
a model response supplies no premises of the proof. Maintainers verify the
namespace assumptions and independent reviewers bind proof, code, and local
gate evidence to the candidate before accepting the claim.

\section*{Key Findings}
The predicate limits the selected pathname to one declared slot. It is
locally complete for eligible, stable, readable slots under the explicit
observation contract. A per-lookup proof composes over a finite trace
without implying that an archive is a transaction or that two run records
belong to distinct principals.

\begin{table}[!htbp]
\centering\small
\caption{Artifact-read results, Part I. Hypotheses are summarized by their dependent result.}
\begin{tabularx}{\linewidth}{@{}p{0.21\linewidth}X l@{}}
\toprule
Case & Result and limit & Source \\
\midrule
Alphabet & One restricted component, subject to host name availability. & \ref{def:inputs} \\
Run registration & Unknown labels are outside the lookup domain. & \ref{def:inputs} \\
Declared stage & Membership precedes filesystem observation. & \ref{def:inputs} \\
Root & Canonical configuration is trusted and fixed for the lookup. & \ref{def:namespace} \\
Expected entry & Exactly four components below the root are prescribed. & \ref{def:namespace} \\
Path aliases & A different canonical slot cannot satisfy equality. & \ref{thm:boundary} \\
\bottomrule
\end{tabularx}
\end{table}
\begin{table}[!htbp]
\centering\small
\caption{Artifact-read results, Part II. Hypotheses are summarized by their dependent result.}
\begin{tabularx}{\linewidth}{@{}p{0.21\linewidth}X l@{}}
\toprule
Case & Result and limit & Source \\
\midrule
Unavailable entry & No selected artifact is produced by a failed observation. & \ref{def:rule} \\
File class & A directory is insufficient. Empty regular bytes can qualify. & \ref{thm:boundary} \\
Case variants & Exact component equality can deny host-equivalent spellings. & \ref{thm:boundary} \\
Handoff & Stable identity through open is a premise of byte selection. & \ref{thm:boundary} \\
Finite trace & Every admitted request satisfies the same local invariant. & \ref{thm:boundary} \\
Archive & Composition limits member names without an atomic ZIP guarantee. & \ref{thm:boundary} \\
\bottomrule
\end{tabularx}
\end{table}
\begin{table}[!htbp]
\centering\small
\caption{Symbols for the exact-slot artifact-read model}
\begin{tabularx}{\linewidth}{@{}lX l@{}}
\toprule
Symbol & Domain and meaning & Source \\
\midrule
$\Sigma$ & Unicode scalar values in a decoded candidate name. & \ref{def:inputs} \\
$A$ & Restricted ASCII alphabet for run labels. & \ref{def:inputs} \\
$D$ & Finite set of registered run labels at admission. & \ref{def:inputs} \\
$S$ & Finite set of declared workflow stages. & \ref{def:inputs} \\
$x$ & Candidate run label, a finite string in $\Sigma^*$. & \ref{def:inputs} \\
$y$ & Candidate stage label, a finite string in $\Sigma^*$. & \ref{def:inputs} \\
$R$ & Canonical trusted directory path. & \ref{def:namespace} \\
$E$ & The unique expected slot for the given pair. & \ref{def:namespace} \\
$C$ & Canonical observation of the proposed entry. & \ref{def:namespace} \\
$\equiv$ & Namespace equality on normalized absolute paths. & \ref{def:namespace} \\
$\denied$ & A distinguished unavailable decision. & \ref{def:rule} \\
$L_i$ & Selection for request $i$ in a finite trace. & \ref{thm:boundary} \\
\bottomrule
\end{tabularx}
\end{table}
\clearpage
\tableofcontents
\listoftables
\clearpage

\section{Scope and source}
The baseline is Engineering Studio AI commit
\nolinkurl{55efcd55a266a629b19ae4926038fdebe5fb0768}. In
\texttt{RunStore.artifact\_path}, a registered run selects a candidate
under the configured storage directory. At that revision the pathname
is returned if it exists. The JSON preview and both download surfaces
consume the same lookup~\cite{runtime}. The run snapshot companion~\cite{snapshot} studies a different
operation, namely publication of \texttt{run.json}; its atomic visibility
result does not establish confinement for artifact reads.
The milestone considered here changes shared read selection only. It
keeps transport schemas and recovery behavior outside the change.

\section{Model and theorem}

\begin{definition}[Inputs and declared names]
\label{def:inputs}
Let $\Sigma^*$ be the set of finite decoded input strings.
Let $A\subset\Sigma$ consist exactly of the ASCII uppercase and
lowercase letters, ASCII decimal digits, underscore, and hyphen.
Let $A^+$ be the nonempty finite strings on $A$.
Let $D\subseteq\Sigma^*$ be the finite run register at admission
and $S\subseteq A^+$ the finite declared stage set.
The pair $(x,y)\in\Sigma^*\times\Sigma^*$ is eligible when
\begin{equation}
\label{eq:eligible}
x\in D\cap A^+,\qquad y\in S.
\end{equation}
Membership is tested on the actual values presented to the shared
lookup. Storage portability additionally depends on the host.
\end{definition}

\begin{definition}[Canonical namespace and slot]
\label{def:namespace}
A path is modeled as a volume or root identifier followed by a
finite sequence of components. Joining a simple component appends
that component to the sequence, without normalizing away any prefix.
Let $R$ be the canonical absolute path of the trusted storage root.
For an eligible pair define the expected path by
\begin{equation}
\label{eq:slot}
E=\slot_R(x,y)=R/x/\texttt{artifacts}/y/\texttt{output.md}.
\end{equation}
The slash here denotes typed component concatenation. The observation
$\resolve(E)$ produces a normalized absolute path $C$ or an unavailable
result $\denied$. A relative configured root is resolved before forming
$E$. Unavailable root resolution also reports $\denied$.
Write $C\equiv E$ for equality of the normalized absolute path component tuples,
including its anchor and case-sensitive string components.
This choice may be stricter than host lookup equality. A regular-file
observation $\regular(C)$ is true only when the platform successfully
classifies $C$ as a regular file.
\end{definition}

\begin{definition}[Selection rule]
\label{def:rule}
A lookup returns $L=\denied$ when any required observation is
unavailable or any gate fails. Otherwise it returns the observed $C$.
Equivalently,
\begin{equation}
\label{eq:rule}
L=
\begin{cases}
C & \text{if \eqref{eq:eligible} holds and $R$ is available,}\\
  & \quad \text{$C=\resolve(E)\ne\denied$, $C\equiv E$, and $\regular(C)$},\\
\denied & \text{otherwise.}
\end{cases}
\end{equation}
An admitted pathname is a certificate of selection under the
hypotheses, subject to a later read succeeding. It does not
contain the result bytes.
\end{definition}

\begin{theorem}[Exact-slot selection and finite composition]
\label{thm:boundary}
Assume the hypotheses below. For one request, if $L\ne\denied$ then
\begin{equation}
\label{eq:claim}
x\in D\cap A^+,\qquad y\in S,\qquad L\equiv\slot_R(x,y).
\end{equation}
The selected entry is a regular file at admission. Any successful
handoff read opens the file at that slot under the observation
premise. Conversely, an eligible pair with an available root,
equal resolved slot, and successful regular-file observation is
admitted. In a finite trace, the implication holds for every
admitted request using its own registry and observation cut.
\end{theorem}

\noindent\textbf{Proof type:} Direct by exhaustive cases, with finite
induction for trace composition.

\noindent\textbf{Hypotheses:}
\begin{enumerate}[label=H\arabic*.,ref=\arabic*]
\item\label{hyp:names} Inputs are decoded strings checked by exact
grammar and exact registry membership before path construction.
Declared stages and the two fixed suffix components join as simple
components on the selected host.
\item\label{hyp:root} $R$ is the canonical path of a trusted storage root.
A configured alias is interpreted relative to that canonical root.
The application controls the directory, and a successful path equality
compares coherent volumes and exact component spellings.
\item\label{hyp:io} Successful resolution and classification report
facts about the same namespace observation. Missing entries, unusable
names, and errors that are handled by the lookup never become positive
classification or equality evidence.
\item\label{hyp:stable} A successful handoff-open selects the file
identified by the admitted observation, or fails. This premise excludes
an identity swap between observation and open. The guarantee concerns
the namespace slot and does not identify the sole name of an inode.
\item\label{hyp:trace} Each request in the finite trace obeys
H\ref{hyp:names}--H\ref{hyp:stable} for its own observations. Each
transport reads only paths returned by the shared boundary.
\end{enumerate}

\noindent\textbf{Proof dependencies:}
\begin{enumerate}
\item Definitions~\ref{def:inputs}--\ref{def:rule} supply the typed
domains, expected path function, and total decision predicate.
\item Exact finite set membership and component concatenation support
the single-request proof. H\ref{hyp:root} supplies the equality premise.
\item H\ref{hyp:stable} connects selection to open. Induction on the
number of processed requests lifts the local result to a finite trace.
\end{enumerate}

\begin{proof}
We prove the local statement by cases over the first failing gate and
the trace statement by induction on the prefix length $n$.

\begin{enumerate}
\item \textbf{Domain and outputs.}
Fix the request label pair $(x,y)$ and the state in which its guards
run. By Definition~\ref{def:inputs} both labels are strings, hence exact
equality and membership in the finite sets are defined. No cardinal
requirement on $D$ is needed. In particular, $D=\varnothing$ is a valid
input state. Although $S$ is finite, its size places no upper bound on
the lengths of the supplied strings. The lookup output is either a
path or the distinguished value $\denied$; these alternatives are
disjoint by their types. We use this distinction before comparing any
paths. If $L=\denied$, \eqref{eq:claim} has a false antecedent.
In that case the theorem makes no assertion that a directory, run,
or stage is absent. Denial also covers failures to obtain sufficiently
reliable observations. If $L\ne\denied$, the rule promises a path,
not yet a descriptor or a byte string. In the cases below, the
first failing gate prevents construction of a later positive result.

\item \textbf{Disjoint exhaustion of the decisions.}
Consider in order stage membership, run spelling, run registration,
availability of root and candidate observations, canonical equality,
and file classification. Each completed Boolean gate has an accepted
or rejected value. Each unavailable observation maps to denial under
H\ref{hyp:io}. Assign an execution to its first failure if there is
one. An execution cannot have two different first failures because
gate positions are totally ordered by their evaluation sequence.
If no gate fails, all gates needed by \eqref{eq:rule} succeed. This
gives the remaining case. The partition remains valid if the file
classification is never reached on a mismatch: a skipped gate has no
value used to overrule the earlier failure. In particular, a directory
found during a denied lookup cannot become an artifact merely because
another component of the request was valid. An infinite environment
failure or a task that never returns is outside the model of completed
lookups; the safety implication is conditional on an actual result.

\item \textbf{Stage boundary.}
Suppose the stage gate fails, so $y\notin S$. Under
H\ref{hyp:names} the gate compares the complete input string against
the fixed stage register. Thus a request can fail even if its run
name is well-formed and registered. It can also fail for an otherwise
existing file. The membership condition in \eqref{eq:eligible} does
not become true from either fact. With $S=\varnothing$, this argument
denies every request. For a singleton set, only the one declared label
can reach later gates. If stage membership succeeds, $S\subseteq A^+$
establishes that the accepted label is nonempty and uses only the
restricted alphabet. The lookup does not infer membership from the
existence of a directory, from output contents, from a substring, or
from a stage label in a model response. Counting admissible stage
labels therefore depends on the workflow declaration and not on the
number of entries visible in the filesystem.

\item \textbf{Run spelling.}
Now assume $y\in S$. If $x$ is empty, then $x\notin A^+$ by the
definition of positive word length. If it is nonempty but contains
any symbol outside $A$, set membership also fails. These two
conditions are exhaustive for failure of spelling. They cover distinct
issues: accepting all-ASCII strings with no length check would
admit the empty word, and accepting every nonempty decoded string
would admit strings with structurally meaningful filesystem symbols.
Conversely, if the length is positive and every character belongs to
$A$, the lexical gate succeeds. This check is on the original decoded
value presented to the lookup. It grants no additional normalization,
Unicode folding, whitespace trimming, or second decoding operation.
A string consisting solely of one allowed character is a legitimate
grammar boundary. Its usability on a host requires separate evidence.
The alphabet excludes volume punctuation, component separators, and
dot components by direct inspection, including their use inside
otherwise ordinary-looking labels. There is no assertion that a very
long word names a valid host entry.

\item \textbf{Run registration.}
If the lexical gate succeeds but $x\notin D$, no expected slot is
selected by the rule. Registration is an equality test on the stored
keys, not a filesystem search. When a directory for that string is
present but the register has no corresponding entry, the request is
still denied. When a registered value fails the grammar, its mere
presence in the register does not make the spelling valid. Thus the
run labels eligible for filesystem checks are the intersection in
\eqref{eq:eligible}. If $D$ is empty, the intersection is empty.
If $D$ has one eligible member, at most that member proceeds.
Multiple observations or multiple stored states for the same key do
not create different strings in the set. This claim treats whichever
registry snapshot the lookup actually observed as the authority.
A later modification to the registry is outside this admission
decision. It does not add a statement about freshness of an already
issued descriptor.

\item \textbf{Expected slot exists as a name.}
Assume both run gates and the stage gate succeed. From the preceding
cases, $x$ and $y$ are nonempty simple lexical components in the
sense of H\ref{hyp:names}. The suffixes \texttt{artifacts} and
\texttt{output.md} are fixed and usable components of the intended
layout, rather than input. Concatenation in \eqref{eq:slot} gives
the unique ordered tuple with prefix $R$ and those four suffix
components. This uniqueness is syntactic. It does not assume an
entry exists at the tuple, nor that the pathname has only one host
spelling. If $R$ has no additional directory components beyond its
volume anchor, suffix concatenation is still defined in the model.
For nested roots or equal strings in distinct volumes, equality keeps
the full root and all components. Thus the canonical-name test cannot
be replaced by a string-prefix check that discards component or
volume identity.

\item \textbf{Available root and namespace stability.}
If the configured root cannot be resolved to the modeled $R$,
\eqref{eq:slot} cannot be used as a certificate and the request is
denied. If the root is available, it is taken from trusted
configuration, never from a run or stage label. In the model a
configuration alias may designate the trusted directory, and the
slots are calculated relative to its canonical form. This allows
administrative placement of the storage root without treating an
arbitrary per-run alias as configuration. The proof holds the choice
fixed while evaluating a request. Later deployment changes to the
working directory, configured root, or filesystem mount require a
new observation. We assume that a successful canonical form denotes
the actual trusted root and does not become a different root before
the selected read. Without this premise a syntactically well-formed
tuple alone could be relative to an untrusted namespace.

\item \textbf{Unavailable candidate.}
For a registered and well-formed pair under the available root,
attempt the candidate observation within the namespace contract.
If this returns $\denied$, the equality condition has no witness.
A missing parent, unavailable volume, or handled I/O failure cannot
supply a canonical path with positive regular-file evidence.
Strict resolution is the chosen implementation contract.
A missing component therefore denies selection at resolution.
For comparison, under a non-strict resolver a partial normalization
could instead reach the final type check. It would still be denied
if the classification cannot return true. The two branches are distinct operationally but
give the same no-selection result. A request with a host-reserved
name illustrates why lexical spelling and name availability are
separate. The alphabet does not decide host device names, path-length
limits, permissions, or share availability. The model requires the
platform to produce coherent absolute-path and regular-file evidence;
a reserved name lacking this evidence cannot be admitted.

\item \textbf{Equality contrast.}
Suppose resolution produces an absolute path $C$. If $C\not\equiv E$,
the failed equality suffices for denial, regardless of the type or
contents of the resolved entry. This covers a resolved entry under
another root, and also a resolved entry under the same root but at a
different declared slot. The second case is the reason for checking
exact slot equality instead of descendant membership alone.
For paths in distinct volumes H\ref{hyp:root} prevents a false
equality merely because remaining components have the same spelling.
For paths in the same volume but with different effective components,
the same premise denies equality. In particular, if an alias occurs
at an intermediate component and canonicalization changes a
significant component, the mismatch is visible without inspecting
payload bytes. This argument uses a positive mismatch and does not
claim that all filesystem redirections are detectable by path
comparison.

\item \textbf{Equality and alias limits.}
If $C\equiv E$, the check compared the actual tuples of
normalized components, including the anchor. Each component
comparison is exact, including case. This comparison is
intentionally stricter than name equality on a host that can
look up several spellings of the same object.
The registry also uses exact string membership. One spelling
might therefore be registered while a second is not, even if
the host accepts both spellings of an emitted pathname.
A lexically valid case variant that resolves to a different
tuple is denied deterministically. If the resolver preserves
both different spellings as reported names, admitting both
does not establish distinct physical ownership. Conversely, if
the resolver canonicalizes both spellings to the same tuple,
at most the one whose expected tuple is equal to it can satisfy
the equality gate. Exact component equality is therefore a
strengthening of the logical slot boundary, but the theorem
does not infer a one-to-one map between admitted pathnames
and physical objects. A deployment requiring such a map
would need a separate storage invariant and evidence.

\item \textbf{Entry classification.}
Suppose all earlier gates succeeded. If $\regular(C)$ is not true,
the final gate denies selection. A directory is not a regular file,
even when its name ends with the expected text. Other nonregular
filesystem types are also excluded by the classification premise.
A permissions problem that yields no positive evidence cannot be
treated as a regular-file observation. If the platform succeeds in
classifying $C$ as regular, size is not considered in this model.
In particular, an empty regular file satisfies this type condition;
it is not equivalent to absence. A regular file with bytes that are
not decodable as the expected text can pass this specific namespace
gate. Payload validation, successful file reads, and successful
transport encoding are separate predicates and require different
error handling. Classifying one pathname also says nothing about
other names for the same physical file.

\item \textbf{Sound selection.}
Consider the final member of the partition, where all gates succeed.
The run spelling and registration gates give $x\in D\cap A^+$ by
intersection introduction. Stage membership gives $y\in S$.
The available root and eligible labels make \eqref{eq:slot} defined.
The candidate observation gives $C=\resolve(E)\ne\denied$ and
the equality gate gives $C\equiv E$.
The classification gives $\regular(C)$.
Using all conditions in \eqref{eq:rule}, we get $L=C$.
Substitution into the canonical equality yields
$L\equiv\slot_R(x,y)$. This is exactly \eqref{eq:claim}.
We also get that the admitted entry is regular at the observation.
The preceding branches all deny selection and hence add no exceptions
to the implication. This proves the single-request selection
property, including for a zero-byte eligible payload and any run
whose declared stage set contains just its selected stage.

\item \textbf{Completeness under the rule.}
For the converse, assume an eligible pair, available root, available
candidate with equal resolved slot, and a successful regular-file
observation for that candidate. Eligibility entails membership in
both run predicates and in $S$ by conjunction elimination.
Thus none of the three input gates fail. The available root allows
construction of the same $E$ from the fixed configuration.
Availability and equality rule out both candidate failure and mismatch.
The regular-file observation rules out the last kind of failure.
The unique remaining partition member returns $C$ by
\eqref{eq:rule}. This is completeness relative to the listed observations
and depends on their consistent collection. It does not claim every
registered run has every stage ready, or every relevant file is
available, or every admitted request is authorized for the particular
caller. A namespace state can satisfy the trust premises while no
positive selection is possible because the store is still empty.

\item \textbf{Read handoff.}
Fix a positive selection $L=C$. By H\ref{hyp:stable}, a successful
open during the intended transport handoff selects the same file
as the admitted observation, or does not hand off a file at all.
If open succeeds, combining this premise with \eqref{eq:claim}
gives the read of the exact slot. If it fails, there is no successful
read to which that conclusion applies. This reasoning is a
conditional refinement, not a claim that obtaining a path pins a
descriptor. An intervening identity replacement can break the
premise even if the pathname strings at both times are equal.
Conversely, an admitted open can continue to use a descriptor that
already denotes the selected object; the proof need not assume
that its bytes stay immutable afterward, beyond what is necessary
for the statement about the selection itself. The present theorem
does not quantify over partial reads or successful parse results.
It is not a model of every system call of the transport.

\item \textbf{Trace base and induction hypothesis.}
Let $T$ be the finite number of completed requests. For each integer
$n$ with $0\le n\le T$, let the prefix property be: every positive
selection among the first $n$ results satisfies \eqref{eq:claim}
using the observations for its own request, and every successful
read in that prefix uses the corresponding selected slot.
For $n=0$ there is no request index in the prefix, so this
universal property holds by the empty-domain rule. No evidence
about the presence of a root directory is entailed by the base case.
If $T=0$, the base case is already the whole proof of trace
composition. Otherwise, assume the prefix property for an arbitrary
$n<T$. The next index $n+1$ is in the finite range.
The $n$-prefix result is unchanged as a statement about past
observations. It does not say that every old decision remains
current after another lookup.

\item \textbf{Trace induction step.}
Apply H\ref{hyp:trace} to request $n+1$. Its local string labels and
registry are now the fixed values of the single-request proof.
If the request is denied, no positive selection is added. If it is
admitted, the single-request soundness conclusion applies.
If there is an associated successful handoff read, the handoff
argument applies under the same request's observation contract.
The trace prefix of length $n+1$ contains exactly the first $n$
requests and that one next request. Combining the induction
hypothesis with the established next case gives the property for
length $n+1$. Both possibilities for its result were exhaustive.
Induction on finite natural numbers proves the property up to
$T$. It does not require requests to be probabilistically
independent. It also does not imply an atomic joint snapshot when
a consumer iterates over a sequence of successful selections.

\item \textbf{Run-level enumeration.}
For a run archive, suppose the caller enumerates declared stage
labels and invokes the shared boundary on each one. Each iteration
is one of the trace requests just proved safe. A denied lookup
supplies no member for that iteration. A successful lookup selects
only its own prescribed slot. If there are no eligible files, the
model permits no included members. If the enumeration uses each
member of $S$ once, the number of selected slots is at most
$|S|$, including the zero and singleton cases. Repeated enumeration of
the same stage can repeat a valid result and is not excluded by the
read-boundary proof. Stable order and archive names can be preserved
separately by keeping the existing stage ordering contract. The
contents may however come from different observation cuts. In
particular, tracking an earlier successful lookup does not make a
later archive operation transactional or make rollback reconstruct
the earlier file set.

\item \textbf{Independence audit.}
No lemma uses a model provider score, elapsed time, a fixed file
size, or evidence about a particular project domain. Positive
membership does not certify the competence of the producing agent.
The proof needs no numerical probability bounds. It reasons about
coherent namespace predicates and stable file identity at the
handoff instead. The registration predicate is not an actor
authorization policy: two distinct callers presenting the same
eligible pair receive the same namespace decision under the same
observations. Nor does the classification premise address hardlinks:
a regular file at one authorized slot can have another physical name
without violating any path equality. Mount mappings can create similar
physical identity questions. Content provenance, integrity checks,
ownership, and isolation of the storage writer can therefore strengthen
a system without being consequences of \eqref{eq:claim}.
The proof also imposes no progress guarantee on I/O that never completes.

\item \textbf{Why weaker premises are insufficient.}
For a formal contrast, take two distinct canonical slots $E$ and
$F$ within the trusted root, with an observation for the expected
entry that selects $F$ as a regular file. Existence and root descent
are both true while equality $F\equiv E$ is false. A rule using only
either weaker fact can admit a wrong slot in this model. For a
different contrast, assume selection observes $E$ at one cut but
a later open follows a different file identity. Even if a prior
canonical equality was true, that earlier fact does not imply
identity of the two opens. This model violates H\ref{hyp:stable}
and shows why it cannot be removed. The examples are logical
systems of observations, not claims about a production
environment. The proof gives no recipe for making either state
occur. A deployment allowing hostile namespace mutation needs
a different, descriptor-level ownership argument.


\item \textbf{Finite candidate and observation counts.}
For an operational scope audit, write $D_A=D\cap A^+$.
Construct the set of eligible request pairs
$U=D_A\times S$. Finite product counting applies because these
are finite sets, and for each distinct first coordinate the pairs
with that coordinate form a disjoint subset indexed by $S$.
Consequently
\begin{equation}
\label{eq:count}
 |U|=|D_A|\,|S|.
\end{equation}
Either zero factor forces zero eligible pairs. If both factors
are positive, there is at least one such pair, irrespective of
whether it has any file. Now define the map from $U$ to
equivalence classes of expected slots by applying
\eqref{eq:slot} and taking the class for $\equiv$.
Its image has cardinality at most $|U|$ by finite image
counting. This bound counts logical candidates, not inode
identifiers, directory entries, open handles, or users.
If the slot map is injective under host equality, each eligible
pair has a different class, so the image cardinality is
exactly \eqref{eq:count}. If any two distinct pairs map to
the same class, a count using distinct names can exceed the
number of slot classes. That possibility is permitted by
H\ref{hyp:root}; a deployment requiring unique physical
ownership would need a stronger invariant. A single request
returns at most one pathname because $R$, $x$, $y$, the suffix,
and $\resolve$ fix the observation used in the rule.
The lookup does not recursively enumerate the filesystem.
This count offers no bound for a filesystem's own resolution
effort, or for read time once a path is selected.

\item \textbf{Declared stage projection and audit identity.}
Let $(x_i,y_i)$, $1\le i\le T$, be the labels actually checked
in the finite trace, and let $D_i$ and $R_i$ denote the
registration snapshot and configured canonical root for that
request. Define the admitted indices by
\begin{equation}
\label{eq:indices}
 I=\{i\in\{1,\ldots,T\}:L_i\ne\denied\}.
\end{equation}
For every $i\in I$, define its audit row to contain the index,
the evaluated run and stage, and the selected slot class.
This is a function on $I$ because each index identifies one
lookup decision and each positive observation names one class.
The theorem entails the row's membership in
$D_i\cap A^+$ and $S$, and entails its equality statement
against $\slot_{R_i}(x_i,y_i)$.
Conversely, a row that records only the class does not
identify the input pair when the slot map is noninjective.
A row that records just the pair does not identify the
slot across configuration-root changes.
These conclusions use only the existence or absence of
preimages under the relevant projections.
For zero admissions, $I=\varnothing$ and the audit row map
has empty domain. No success should be inferred from an empty
audit table. Repeated requests at different indices give distinct
observations even when their admitted class is the same.
The trace model supplies a reason to bind audit evidence to
observations and candidate versions without requiring the
application to emit raw paths or sensitive artifact bytes.

\item \textbf{Temporal comparison of distinct decisions.}
Take two completed requests with identical evaluated labels and
run registration. Suppose that the canonical root and the
resolved path are equivalent under $\equiv$ across the two
observations, and both classify the slot as regular.
Applying completeness to each request produces two admissions.
By the transitivity of path equivalence, they select the same
logical class of slot. This conclusion does not establish
that the corresponding physical object or bytes were the same:
H\ref{hyp:stable} speaks within a request, and has no relation
between the two observations.
Now hold labels, register, and root fixed but change the
availability of the file between requests.
A regular-file observation at one cut and an unavailable
observation at the other can produce one admission and one denial.
Neither result contradicts the trace statement because the
proof uses a different observation for each decision.
Thus caching a positive result indefinitely would not refine
the model unless the required observation stays valid.
For the opposite ordering, a denied lookup followed by a
regular-file observation allows the second decision to admit.
It does not mean the denial was incorrect when issued.
Availability is a property of the request and namespace state
together; stability of a string pair alone does not extend it.

\item \textbf{Scope of monotonicity.}
Suppose two registrations satisfy $D_1\subseteq D_2$.
If $x\in D_1\cap A^+$, then $x\in D_2\cap A^+$, by the
definition of subset and intersection. With all other
observations held fixed and coherent, every earlier
admission therefore remains admitted under the enlarged
register. The same implication holds when the declared
stage set is extended and an old stage remains unchanged.
In each argument, \eqref{eq:rule} justifies reusing the
same root, candidate, equality, and type witnesses.
This is a restricted monotonicity result about sets
with fixed filesystem observations. A namespace change
can remove or invalidate an observation regardless of
registry inclusion. Also, shrinking a register can
cause rejection without changing any bytes.
If a release is reverted, the old code need not compute
the enlarged or strengthened predicate. No time-travel
claim follows: a previously successful return remains
part of the completed trace, even after policy changes.
Rollback planning therefore needs a distinct release
identity and an explicit statement about which earlier
read decisions cannot be compensated.

\item \textbf{Why a guard is needed at the shared boundary.}
The model has no premise that all callers use the same
frontend. Instead, H\ref{hyp:trace} requires each consumer
to obtain its read path through the shared rule. Consider
a family of consumer functions $f_j$ indexed by $j\in J$,
where each function's domain is a completed selection
and the accompanying request context.
If a consumer receives $\denied$, it has no selected
path to open; its response policy can map the result to
an existing unavailable case or omit a collection member.
If it receives a path, the assumption on consumers
requires any successful selected open to use that path.
For each $j$, the single-request result therefore applies
independently of the presentation of the output.
If $J$ is empty, this universal statement is vacuous.
If an additional consumer independently constructs a
different path, the hypothesis no longer holds for that
consumer and the theorem supplies no safety conclusion.
A bounded review can check the three concrete call sites
even without making any promise about future tools.
This argument separates a single library boundary from
the transport contract: identical unavailable outcomes
can intentionally hide the exact failed predicate.

\item \textbf{Comparison of expected and observed spellings.}
For clarity, suppose volume and normalized component representation
are compared by the same exact tuple operation for both paths.
For a given request the expected tuple has length $m+4$ when
$m$ counts the root components in the chosen representation.
The tuple equality test has only two cases. If the lengths
differ, the tuples cannot be equal, regardless of the agreement
of their shorter prefix. If the lengths agree, equality requires
pointwise equality at every index. In this case an unequal
volume counts as a differing anchor in the normalized
representation; it cannot be hidden in the suffix.
Likewise one differing run spelling, suffix literal, or stage
spelling makes the equality false. A match on every index
constructs the positive witness. This finite comparison does not
rely on a case-folding relation. Host spellings that compare
equal under a separate case-folding rule do not satisfy the test
unless these exact components match. Names that are lexically
distinct at an index can still alias an object below the
comparison layer. In that event the theorem remains a statement
about exact canonical observations and their selected names.
No filesystem case policy has to be modified to apply the check.
For the degenerate root representation with no trailing directory
components, the four appended components still exist and volume
comparison is still part of the full path comparison.
Thus the exact check does not reduce to checking just the
basename, final suffix, or the first shared path component.

\item \textbf{Auditor checklist and limits.}
A reproducibility record should identify the exact source revision,
the model version, and the expected stage sequence. Check that the
serialized request values and stage allowlist at the central lookup
are the actual inputs to the grammar and membership predicates.
Check that root canonicalization precedes slot construction and
that equality and regular-file checks gate every positive return.
Check that preview, single download, and archive consumers obtain
their selected path only through that boundary. Confirm that
rejection uses the existing missing-artifact path and does not leak
storage locations. Report observations of unavailable aliases and
platform features distinctly from clean tests. Verify that no
directory or storage protocol under review invalidates
H\ref{hyp:root} or H\ref{hyp:stable}. Record the test and resource
bounds separately from the model: finite trace composition by itself
proves neither bounded response time nor storage consumption.
Independent review and local checks can verify a candidate against
these obligations but cannot infer untested platform assumptions.
\qedhere
\end{enumerate}
\end{proof}

\section{Refinement and regression obligations}
\label{sec:refinement}
The inspectable refinement uses strict canonical resolution and exact comparison of the path component tuples. It is one defensive gate in
\texttt{RunStore.artifact\_path}. Keep the established run registration
check and stage vocabulary, and compare the candidate's canonical
observation with its expected absolute slot. A returned canonical
path denotes the same configured slot as the ordinary output layout.
An I/O error later in a transport read is not translated by the lookup. File type and handled lookup I/O failures must be tested along with valid
selection. Verify ordinary ready and missing artifacts, rejected
selectors, disallowed or redirected slots, files in distinct
known stages, and transport responses that depend on selection.
This document makes a conditional safety claim about one namespace
decision. The actual platform evidence must record which tests use
real temporary paths, which simulate filesystem errors, and which
cannot exercise alias behavior because of host privileges. Retain
local source and dependency scanning results with the exact revision.
An unavailable or stale advisory source cannot support an assertion
that a dependency has no known advisory. The release record should
bind the proof, regression evidence, and prior accepted revision so
the source change can be reverted without silently rewriting run
state or claiming that earlier disclosures can be undone.

\begin{thebibliography}{9}
\bibitem{runtime}
Engineering Studio AI source at commit
\nolinkurl{55efcd55a266a629b19ae4926038fdebe5fb0768}.
See \texttt{src/engineering\_studio/runs.py} and the run and
download API modules for the observed lookup and consumers.
\bibitem{snapshot}
Engineering Studio AI, \textit{Atomic Publication of Engineering
Studio Run Snapshots}, OpenAI Dev Week companion, 2026.
This companion supplies the distinct write-visibility model for
the run history sidecar.
\end{thebibliography}

\clearpage
\section*{Changelog}
\addcontentsline{toc}{section}{Changelog}
\begin{table}[!htbp]
\centering\small
\caption{Revision history}
\begin{tabularx}{\linewidth}{@{}llX@{}}
\toprule
Version & Date & Change \\
\midrule
2026.1.0.0 & 2026-11-20 & Introduced exact-slot read selection, limits
of namespace evidence, and finite composition across artifact consumers. \\
\bottomrule
\end{tabularx}
\end{table}
\end{document}
